\chapter*{Conclusion}
\addcontentsline{toc}{chapter}{Conclusion}
\markboth{Conclusion}{Conclusion}
\label{chap:conclusion}
\setlength{\parskip}{1em}

The management of supply chain security within critical IT infrastructure remains a vital concern, as the frequency of incidents continues to rise despite the development of new standards and tools. Given the current instability of the digital landscape, a natural shift toward preventive measures during the initial stages of software production and distribution is generally expected. However, our initial simulation iteration is insufficient to fully demonstrate these trends. A more detailed and repetitive approach is necessary to uncover deeper patterns and analyze real-world supply chain attacks more comprehensively. Specifically, leveraging Graph Neural Networks (GNNs) to analyze microservices connections and package dependencies might allow for a more robust identification of patterns and connections within the ecosystem.

Repurposing legacy smartphones to build decentralized networks during hardware shortages remains a critical goal, though many older devices risk becoming electronic waste due to poor performance and a lack of crucial security updates. The possible success of this initiative depends heavily on its technical execution and whether it can offer a genuinelyuseful, viable, and feasible alternative to current centralized systems. At this stage, the project serves as an experimental foundational draft, meaning its precise operational mechanisms have not yet been fully defined. These complex system dynamics and actor interactions require ongoing, iterative research and refinement well beyond the scope of a single introductory article. Furthermore, significant technical hurdles remain unresolved, particularly concerning strict Android security policies that often require device verification and block the execution of uncertified applications.

Future works will focus on the full practical implementation of the proposed system, alongside a rigorous analysis of software components to validate the conceptual framework using real-world data and gathered statistics. We are planning a more advanced experiment utilizing an adversarial simulation that leverages local language models to represent users more accurately and dynamically. Each agent will be provided with a unique persona that influences their actions and social interactions, such as discussing news or reporting issues within community groups. By hosting this within an environment of multiple virtual machines and services, we can create a much more realistic representation of reality compared to simulations based purely on pseudo-random generation. However, the integration of AI agents as “Interactors” currently presents reliability challenges, specifically regarding vulnerabilities like cross-prompt injection, as they require creating and adding guardrails.

Also, we will prioritize essential auxiliary infrastructure, specifically the construction of custom- engineered device farms using modular, 3D-printed server racks equipped with active cooling and continuous power to prevent overheating and rapid battery depletion. Additionally, we plan to systematically gather detailed data across various low-end device categories to establish precise performance benchmarks, hardware characteristics, and application compatibility profiles. There is also a significant opportunity to enhance the network's AI capabilities by implementing self- configuring, distributed inference loops and meta-prompting via tools like llama-cpp and llama-bench, optimized specifically for low-resource hardware. Beyond our primary system architecture, these repurposed smartphones can support broader practical applications by leveraging Termux, proot-distro for Debian packages, and Android Debug Bridge (ADB) connections. However, successfully porting such applications to this ecosystem remains a notable challenge, largely due to the inherent technical difficulties of running systemd natively on Android, even when utilizing workarounds like termux-services. For example, we consider different applications:
\begin{itemize}
    \item Custom inference servers by using the llama-cpp engine to host localized language models, utilizing RPC to distribute computational workloads across the device network.
    \item Decentralized storage and hosting by integrating IPFS network for persistent data storage of open-source repositories, alongside distributed web hosting that serves resources via individual devices using a load balancer.
    \item General distributed computing for execution of general-purpose tasks, including cryptocurrency mining operations using specialized software like cpuminer connected through a stratum server.
\end{itemize}
